
Most machine learning datasets are invisible --- not because they lack quality, but because they lack tags. On OpenML, a platform hosting more than 5,000 actively-studied datasets, datasets with at least one keyword tag register 22.6\% more ML task experiments than their untagged counterparts (IRR=1.2263, 95\% CI [1.1681, 1.2873], p=1.$87\times 10^{-16}$, N=5,217), after controlling for dataset size, age, and decade of upload. This finding is statistically robust across seven model variants, and reflects a structural platform property rather than a temporal artifact.

The disparity in ML dataset adoption is a well-documented but incompletely understood phenomenon. Some datasets accumulate hundreds of task registrations; others, of comparable size and scope, remain largely unused. Explanations typically focus on intrinsic dataset properties --- instance count, dimensionality, task difficulty --- as primary drivers of engagement. These structural features are controlled for explicitly in the present study. They cannot, however, account for adoption gaps among datasets of similar size and complexity. A complementary mechanism is required.

The mechanism examined here is discoverability via search-indexed metadata. On platforms such as OpenML, researchers discover datasets primarily through keyword search. A dataset without keyword tags is structurally absent from tag-indexed search results, regardless of its intrinsic merit. This is the FAIR F1 (Findability) principle operationalized: keyword tags are the mechanism by which datasets enter researcher awareness. The FAIR Data Principles \citep{Wilkinson2016FAIR} established this as a foundational requirement for dataset reuse, but quantitative IRR evidence for the effect of FAIR F1 compliance on ML dataset adoption has been absent from the literature.

The FAIR F1 tagging mechanism on OpenML operates as a \textbf{threshold-plus-amplification} structure. Binary tag presence (has\_tags=1 vs. has\_tags=0) captures the critical search-index-membership threshold: any tagging is associated with a 22.6\% adoption advantage. Above this threshold, tag count magnitude amplifies discovery probability log-linearly: within the tagged subset (N=2,625), each unit increase in log(tag\_count+1) is associated with 53.3\% more task registrations (IRR=1.5332, p=1.$28\times 10^{-82})$. The categorical shape of this amplification is non-uniform --- the 6+ tag tier produces a meaningfully distinct effect (IRR=1.2861) while 1--5 tags produce statistically indistinguishable effects --- indicating that comprehensive tagging maximizes discovery benefit.

The study makes four contributions:

\textbf{C1.} First empirical NB-2 quantification of FAIR F1 keyword tagging and ML dataset adoption. Incidence rate ratio estimates are provided for has\_tags binary (IRR=1.2263) and log tag count (IRR=1.5332) in negative binomial regression appropriate for overdispersed count outcomes, on the OpenML platform with N=5,217 datasets.

\textbf{C2.} Threshold-plus-amplification characterization of the FAIR F1 mechanism. The mechanism shape is characterized as binary threshold (any tags $\rightarrow$ search index membership), continuous log-linear amplification (more tags $\rightarrow$ more search pathways), and non-uniform categorical tiers (6+ tag dominance).

\textbf{C3.} Structural mechanism verification via decade fixed effects. The has\_tags effect survives C(decade) fixed effects with 12.2\% attenuation (versus 98.6\% attenuation for a prior composite metadata score tested on the same corpus), establishing the within-decade, structural character of the finding.

\textbf{C4.} Informative negative for uniform categorical dose-response. Categorical analysis reveals that OpenML tagging behavior is bimodal --- users tag comprehensively (3+ tags) or not at all --- making the 1--2 tag range a sparse middle ground (N=73). This constrains future tagging intervention designs.

Section 2 reviews the FAIR principles literature, ML metadata adoption research, and negative binomial regression methodology. Section 3 describes the study design, corpus, and model specification. Section 4 presents the four experimental analyses. Section 5 reports results. Section 6 interprets findings and acknowledges limitations. Section 7 concludes.

